\section{Methodology}
\label{sec:methodology}

Our model is developed in four phases of increasing complexity, each adding
one feature to the network. We begin with a single gas whose reference
reading is the parent of all five sensors, to test how well the sensors fit
the data. We then add the remaining gases and cross-sensitivity edges,
learning the CPTs with maximum likelihood. Then, we make the true gas levels
latent, learn them with EM~\citep{dempster1977em}, and compute hourly air
quality as the worst band across gases; finally, as a reach goal, we extend
the model to an HMM over hours, learned with Baum--Welch~\citep{rabiner1989hmm}.
In each phase, we use exact inference to compute the distribution over a gas's
level given the sensor readings and weather. To evaluate, we train on one block
of time and test on a later, held-out block: on test hours with a reference
reading, we hide it, infer it from the sensors, and report accuracy and
log-likelihood against a baseline that predicts from the hour of day alone.
We also compare configurations, varying the number of discretization levels
and including or excluding the cross-sensitivity edges.
